%% build.tex
%%
%% Copyright (C) 1999-2018 Simone Piccardi.  Permission is granted to copy,
%% distribute and/or modify this document under the terms of the GNU Free
%% Documentation License, Version 1.1 or any later version published by the
%% Free Software Foundation; with the Invariant Sections being "Un preambolo",
%% with no Front-Cover Texts, and with no Back-Cover Texts.  A copy of the
%% license is included in the section entitled "GNU Free Documentation
%% License".
%%

\chapter{Gli strumenti di ausilio per la programmazione}
\label{cha:build_manage}

Tratteremo in questa appendice in maniera superficiale i principali strumenti
che vengono utilizzati per programmare in ambito Linux, ed in particolare gli
strumenti per la compilazione e la costruzione di programmi e librerie, e gli
strumenti di gestione dei sorgenti e di controllo di versione.

Questo materiale è ripreso da un vecchio articolo, ed al momento è molto
obsoleto.


\section{L'uso di \texttt{make} per l'automazione della compilazione}
\label{sec:build_make}

Il comando \texttt{make} serve per automatizzare il processo di costruzione di
un programma ed effettuare una compilazione intelligente di tutti i file
relativi ad un progetto software, ricompilando solo i file necessari ed
eseguendo automaticamente tutte le operazioni che possono essere necessarie
alla produzione del risultato finale.\footnote{in realtà \texttt{make} non si
  applica solo ai programmi, ma in generale alla automazione di processi di
  costruzione, ad esempio anche la creazione dei file di questa guida viene
  fatta con \texttt{make}.}


\subsection {Introduzione a \texttt{make}}
\label{sec:make_intro}

Con \texttt{make} si possono definire i simboli del preprocessore C che
consentono la compilazione condizionale dei programmi (anche in Fortran); è
pertanto possibile gestire la ricompilazione dei programmi con diverse
configurazioni con la modifica di un unico file.

La sintassi normale del comando (quella che si usa quasi sempre, per le
opzioni vedere la pagina di manuale) è semplicemente \texttt{make}. Questo
comando esegue le istruzioni contenute in un file standard (usualmente
\texttt{Makefile}, o \texttt{makefile} nella directory corrente).

Il formato normale dei comandi contenuti in un \texttt{Makefile} è:
\begin{verbatim}
bersaglio: dipendenza1 dipendenza2 ...
        regola1
        regola2
        ...
\end{verbatim}
dove lo spazio all'inizio deve essere un tabulatore (metterci degli spazi è un
errore comune, fortunatamente ben segnalato dalle ultime versioni del
programma), il bersaglio e le dipendenze nomi di file e le regole comandi di
shell.

Il concetto di base è che se uno dei file di dipendenza è più recente (nel
senso di tempo di ultima modifica) del file bersaglio quest'ultimo viene
ricostruito di nuovo usando le regole elencate nelle righe successive.

Il comando \texttt{make} ricostruisce di default il primo bersaglio che viene
trovato nella scansione del \texttt{Makefile}, se in un \texttt{Makefile} sono
contenuti più bersagli indipendenti, si può farne ricostruire un altro che non
sia il primo passandolo esplicitamente al comando come argomento, con qualcosa
del tipo di: \texttt{make altrobersaglio}.

Si tenga presente che le dipendenze stesse possono essere dichiarate come 
bersagli dipendenti da altri file; in questo modo è possibile creare una 
catena di ricostruzioni. 

In esempio comune di quello che si fa è mettere come primo bersaglio il
programma principale che si vuole usare, e come dipendenze tutte gli oggetti
delle funzioni subordinate che utilizza, con i quali deve essere collegato; a
loro volta questi oggetti sono bersagli che hanno come dipendenza i relativi
sorgenti. In questo modo il cambiamento di una delle funzioni subordinate
comporta solo la ricompilazione della medesima e del programma finale.


\subsection{Utilizzo di \texttt{make}}

Il comando \texttt{make} mette a disposizione una serie molto complesse di
opzioni e di regole standard predefinite e sottintese, che permettono una
gestione estremamente rapida e concisa di progetti anche molto complessi; per
questo piuttosto che fare una replica del manuale preferisco commentare un
esempio di \texttt{makefile}, quello usato per ricompilare i programmi di
analisi dei dati dei test su fascio del tracciatore di Pamela.

\small
\begin{verbatim}
#----------------------------------------------------------------------
#
# Makefile for a Linux System:
# use GNU FORTRAN compiler g77
# Makefile done for tracker test data
#
#----------------------------------------------------------------------
# Fortran flags
FC=g77
FFLAGS= -fvxt -fno-automatic -Wall -O6 -DPC # -DDEBUG
CC=gcc
CFLAGS= -Wall -O6
CFLADJ=-c #-DDEBUG
#
# FC            Fortran compiler for standard rules
# FFLAGS        Fortran flags for standard rules
# CC            C Compiler for standard rules
# CFLAGS        C compiler flags for standard rules 
LIBS=  -L/cern/pro/lib -lkernlib -lpacklib -lgraflib -lmathlib
OBJ=cnoise.o fit2.o pedsig.o loop.o badstrp.o cutcn.o readevnt.o \
erasepedvar.o readinit.o dumpval.o writeinit.o

riduzione: riduzione.F $(OBJ) commondef.f readfile.o
        $(FC) $(FFLAGS) -o riduzione riduzione.F readfile.o $(OBJ) $(LIBS) 

readfile.o: readfile.c
        $(CC) $(CFLAGS) -o readfile.o readfile.c

$(OBJ): commondef.f

.PHONY : clean
clean:
        rm -f *.o 
        rm -f *~
        rm -f riduzione
        rm -f *.rz
        rm -f output
\end{verbatim}
\normalsize

Anzitutto i commenti, ogni linea che inizia con un \texttt{\#} è un
commento e non viene presa in considerazione.

Con \texttt{make} possono essere definite delle variabili, da potersi riusare
a piacimento, per leggibilità si tende a definirle tutte maiuscole,
nell'esempio ne sono definite varie: \small
\begin{verbatim}
FC=g77
FFLAGS= -fvxt -fno-automatic -Wall -O6 -DPC # -DDEBUG
CC=gcc
CFLAGS= -Wall -O6
CFLADJ=-c #-DDEBUG
...
LIBS=  -L/cern/pro/lib -lkernlib -lpacklib -lgraflib -lmathlib
OBJ=cnoise.o fit2.o pedsig.o loop.o badstrp.o cutcn.o readevnt.o \
\end{verbatim}
\normalsize

La sintassi è \texttt{NOME=}, alcuni nomi però hanno un significato speciale
(nel caso \texttt{FC}, \texttt{FLAGS}, \texttt{CC}, \texttt{CFLAGS}) in quanto
sono usati da \texttt{make} nelle cosiddette \textsl{regole implicite} (su cui
torneremo dopo).

Nel caso specifico, vedi anche i commenti, abbiamo definito i comandi di 
compilazione da usare per il C e il Fortran, e i rispettivi flag, una variabile
che contiene il pathname e la lista delle librerie del CERN e una variabile con una
lista di file oggetto.

Per richiamare una variabile si usa la sintassi \texttt{\$(NOME)}, ad esempio
nel makefile abbiamo usato:
\begin{verbatim}
        $(FC) $(FFLAGS) -o riduzione riduzione.F readfile.o $(OBJ) $(LIBS) 
\end{verbatim}
e questo significa che la regola verrà trattata come se avessimo scritto 
esplicitamente i valori delle variabili.

Veniamo ora alla parte principale del makefile che esegue la costruzione del
programma:
\begin{verbatim}
riduzione: riduzione.F $(OBJ) commondef.f readfile.o
        $(FC) $(FFLAGS) -o riduzione riduzione.F readfile.o $(OBJ) $(LIBS) 
readfile.o: readfile.c
        $(CC) $(CFLAGS) -o readfile.o readfile.c
$(OBJ): commondef.f
\end{verbatim}

Il primo bersaglio del makefile, che definisce il bersaglio di default, è il
programma di riduzione dei dati; esso dipende dal suo sorgente da tutti gli
oggetti definiti dalla variabile \texttt{OBJ}, dal file di definizioni
\texttt{commondef.f} e dalla routine C \texttt{readfile.o}; si noti il
\texttt{.F} del sorgente, che significa che il file prima di essere compilato
viene fatto passare attraverso il preprocessore C (cosa che non avviene per i
\texttt{.f}) che permette di usare i comandi di compilazione condizionale del
preprocessore C con la relativa sintassi. Sotto segue il comando di
compilazione che sfrutta le variabili definite in precedenza per specificare
quale compilatore e opzioni usare e specifica di nuovo gli oggetti e le
librerie.

Il secondo bersaglio definisce le regole per la compilazione della routine
in C; essa dipende solo dal suo sorgente. Si noti che per la compilazione
vengono usate le variabili relative al compilatore C. Si noti anche che se
questa regola viene usata, allora lo sarà anche la precedente, dato che
\texttt{riduzione} dipende da \texttt{readfile.o}.

Il terzo bersaglio è apparentemente incomprensibile dato che vi compare
solo il riferimento alla variabile \texttt{OBJ} con una sola dipendenza e
nessuna regola, essa però mostra le possibilità (oltre che la
complessità) di \texttt{make} connesse alla presenza di quelle regole
implicite a cui avevamo accennato.

Anzitutto una peculiarità di \texttt{make} è che si possono anche usare più
bersagli per una stessa regola (nell'esempio quelli contenuti nella variabile
\texttt{OBJ} che viene espansa in una lista); in questo caso la regola di
costruzione sarà applicata a ciascuno che si potrà citare nella regola stessa
facendo riferimento con la variabile automatica: \texttt{\$@}.  L'esempio
usato per la nostra costruzione però sembra non avere neanche la regola di
costruzione.

Questa mancanza sia di regola che di dipendenze (ad esempio dai vari sorgenti) 
illustra le capacità di funzionamento automatico di \texttt{make}.
Infatti è facile immaginarsi che un oggetto dipenda da un sorgente, e che 
per ottenere l'oggetto si debba compilare quest'ultimo.

Il comando \texttt{make} sa tutto questo per cui quando un bersaglio è un
oggetto (cioè ha un nome tipo \texttt{qualcosa.o}) non è necessario
specificare il sorgente, ma il programma lo va a cercare nella directory
corrente (ma è possibile pure dirgli di cercarlo altrove, il caso è trattato
nel manuale).  Nel caso specifico allora si è messo come dipendenza solo il
file delle definizioni che viene incluso in ogni subroutine.

Inoltre come dicevamo in genere per costruire un oggetto si deve compilarne il
sorgente; \texttt{make} sa anche questo e sulla base dell'estensione del
sorgente trovato (che nel caso sarà un \texttt{qualcosa.f}) applica la regola
implicita.  In questo caso la regola è quella di chiamare il compilatore
fortran applicato al file oggetto e al relativo sorgente, questo viene fatto
usando la variabile \texttt{FC} che è una delle variabili standard usata dalle
regole implicite (come \texttt{CC} nel caso di file \texttt{.c}); per una
maggiore flessibilità poi la regola standard usa anche la variabile
\texttt{FFLAGS} per specificare, a scelta dell'utente che non ha che da
definirla, quali flag di compilazione usare (nella documentazione sono
riportate tutte le regole implicite e le relative variabili usate).

In questo modo è stato possibile usare una sola riga per indicare la serie di
dipendenze e relative compilazioni delle singole subroutine; inoltre con l'uso
della variabile \texttt{OBJ} l'aggiunta di una nuova eventuale routine
\texttt{nuova.f} comporta solo l'aggiunta di \texttt{nuova.o} alla definizione
di \texttt{OBJ}.


\section{Source Control Management}
\label{sec:build_scm}

Uno dei problemi più comuni che si hanno nella programmazione è quella di
poter disporre di un sistema che consenta di tenere conto del lavoro
effettuato, di tracciare l'evoluzione del codice, e, soprattutto nel caso di
progetti portati avanti da più persone, consentire un accesso opportunamente
coordinato fra i vari partecipanti alla base comune dei sorgenti dello
sviluppo.

I programmi che servono a questo scopo vanno sotto il nome comune di SCM
(\textit{Source Control Manager}), e ne esistono di diversi tipi con diverse
filosofie progettuali, in particolare nelle modalità con cui gestiscono
l'accesso alla base di codice comune da parte dei singoli programmatori che vi
accedono.

Fra questi uno dei più usati, nonostante la sua architettura sia considerata
superata, è Subversion, un sistema di archiviazione centralizzata del codice
che consente di tenere traccia di tutte le modifiche e di condividere un
archivio comune per progetti portati avanti da diverse persone.

\subsection{Introduzione a Subversion}
\label{sec:build_subversion}

Subversion è basato sul concetto di \textit{repository}, un archivio
centralizzato in cui vengono riposti e da cui vengono presi i sorgenti dei
programmi. L'archivio tiene traccia delle diverse versioni registrate; i
programmatori inviano le modifiche usando una copia locale che hanno nella
loro directory di lavoro.

Subversion può gestire più di un progetto all'interno di un singolo server,
ciascuno dei quali viene associato ad un \textit{repository} distinto, ma si
possono anche creare sotto-progetti suddividendo un \textit{repository} in
diverse directory; ma ciascun progetto avrà meccanismi di controllo (ad 
esempio quelli che consentono di inviare email all'inserimento di nuovo
codice) comuni.

Una delle caratteristiche che contraddistinguono Subversion dal suo
predecessore CVS è quella di essere gestibile in maniera molto flessibile
l'accesso al repository, che può avvenire sia in maniera diretta facendo
riferimento alla directory in cui questo è stato installato che via rete,
tramite diversi protocolli. L'accesso più comune è fatto direttamente via
HTTP, utilizzando opportune estensioni del protocollo DAV, ma è possibile
passare attraverso SSH o fornire un servizio di rete dedicato.\footnote{esiste
  all'uopo il programma \texttt{svnserve}, ma il suo uso è sconsigliato per le
  scarse prestazioni e le difficoltà riscontrate a gestire accessi di utenti
  diversi; la modalità di accesso preferita resta quella tramite le estensioni
  al protocollo DAV.}

In generale è comunque necessario preoccuparsi delle modalità di accesso al
codice soltanto in fase di primo accesso al \textit{repository}, che occorrerà
identificare o con il pathname alla directory dove questo si trova o con una
opportuna URL (con il comune accesso via web del tutto analoga a quella che si
usa in un browser), dopo di che detto indirizzo sarà salvato nella propria
copia locale dei dati ed il riferimento diventerà implicito.

Il programma prevede infatti che in ogni directory che si è ottenuta come
copia locale sia presente una directory \texttt{.svn} contenente tutti i dati
necessari al programma. Inoltre il programma usa la directory
\texttt{.subversion} nella home dell'utente per mantenere le configurazioni
generali del client e le eventuali informazioni di autenticazione. 

Tutte le operazioni di lavoro sul \textit{repository} vengono effettuate lato
client tramite il comando \texttt{svn} che vedremo in
sez.~\ref{sec:build_subversion} ma la creazione e la inizializzazione dello
stesso (così come la gestione lato server) devono essere fatte tramite il
comando \texttt{svnadmin} eseguito sulla macchina che lo ospita. In generale
infatti il comando \texttt{svn} richiede che si faccia riferimento ad un
\textit{repository} (al limite anche vuoto) esistente e questo deve essere
opportunamente creato.

Il comando \texttt{svnadmin} utilizza una sintassi che richiede sempre
l'ulteriore specificazione di un sotto-comando, seguito da eventuali altri
argomenti. L'inizializzazione di un \textit{repository} (che sarà creato
sempre vuoto) viene eseguita con il comando:
\begin{verbatim}
svnadmin create /path/to/repository
\end{verbatim}
dove \texttt{/path/to/repository} è la directory dove verranno creati e
mantenuti tutti i file, una volta creato il \textit{repository} si potrà
iniziare ad utilizzarlo ed inserirvi i contenuti con il comando \texttt{svn}. 

Non essendo questo un testo di amministrazione di sistema non tratteremo qui i
dettagli della configurazione del server per l'accesso via rete al
\textit{repository}, per i quali si rimanda alla documentazione del progetto
ed alla documentazione sistemistica scritta per Truelite
Srl.\footnote{rispettivamente disponibili su \url{svn.tigris.org} e
  \url{labs.truelite.it/truedoc}.}

\subsection{Utilizzo di \texttt{svn}}
\label{sec:build_svn_use}

Una volta che si abbia a disposizione un \textit{repository} si potrà creare
un nuovo progetto sottoposto a controllo di versione importando al suo interno
i dati disponibili. In genere è pratica comune suddividere il contenuto di un
repository in tre directory secondo il seguente schema:
\begin{basedescript}{\desclabelwidth{2.7cm}\desclabelstyle{\nextlinelabel}}
\item[\texttt{trunk}] contiene la versione corrente i sviluppo, su cui vengono
  effettuate normalmente le modifiche e gli aggiornamenti;
\item[\texttt{tags}] contiene le diverse versioni \textsl{fotografate} ad un
  certo istante del processo di sviluppo, ad esempio in occasione del rilascio
  di una versione stabile, così che sia possibile identificarle facilmente;
\item[\texttt{branches}] contiene \textsl{rami} alternativi di sviluppo, ad
  esempio quello delle correzioni eseguite ad una versione stabile, che
  vengono portati avanti in maniera indipendente dalla versione principale.
\end{basedescript}

Questa suddivisione consente di sfruttare la capacità di Subversion di creare
senza spesa copie diverse del proprio contenuto, pertanto in genere si pone il
proprio progetto di sviluppo sotto \texttt{trunk}, e si copia quest'ultima in
occasione delle varie versioni di rilascio in altrettante sottocartelle di
\texttt{tags} e qualora si voglia aprire un ramo alternativo di sviluppo
basterà copiarsi il punto di partenza del ramo sotto \texttt{branches} e
iniziare ad eseguire le modifiche su di esso.

Le operazioni di gestione di un progetto con Subversion vengono eseguite con
il comando \texttt{svn}, che analogamente al precedente \texttt{svnadmin}
utilizza una sintassi basata sulla specificazione degli opportuni
sotto-comandi. Si sono riportati quelli più importanti in
tab.~\ref{tab:svn_mainsubcommands}. 

\begin{table}[htb]
  \centering
  \footnotesize
  \begin{tabular}[c]{|l|l|p{8cm}|}
    \hline
     \multicolumn{2}{|c|}{\textbf{Sotto-comando}}& \textbf{Significato} \\
    \hline
    \hline
    \texttt{import}  &  --       & Importa i file della directory corrente sul
                                   \textit{repository}.\\ 
    \texttt{checkout}&\texttt{co}& Scarica una versione del progetto dal
                                   \textit{repository}.\\
    \texttt{commit}  &\texttt{ci}& Invia le modifiche effettuate localmente al 
                                   \textit{repository}.\\
    \texttt{add}     &  --       & Richiede l'aggiunta un file o una directory
                                   al \textit{repository}.\\
    \texttt{remove}  &\texttt{rm}& Richiede la rimozione un file o una 
                                   directory dal \textit{repository}.\\
    \texttt{copy}    &\texttt{cp}& Richiede la copia un file o una cartella del
                                   progetto (mantenendone la storia).\\
    \texttt{move}    &\texttt{mv}& Richiede lo spostamento un file o una 
                                   directory (equivalente ad un \texttt{cp}
                                   seguito da un \texttt{rm}).\\
    \texttt{update}  &  --       & Aggiorna la copia locale.\\
    \texttt{resolved}&  --       & Rimuove una situazione di conflitto presente
                                   su un file.\\ 
    \hline
  \end{tabular}
  \caption{Tabella riassuntiva dei principali sotto-comandi di \texttt{svn}.}
  \label{tab:svn_mainsubcommands}
\end{table}

In genere però è piuttosto raro iniziare un progetto totalmente da zero, è
molto più comune avere una qualche versione iniziale dei propri file
all'interno di una cartella. In questo caso il primo passo è quello di
eseguire una inizializzazione del \textit{repository} importando al suo
interno quanto già esistente.  Per far questo occorre eseguire il comando:
\begin{verbatim}
svn import [/pathname] URL
\end{verbatim}
questo può essere eseguito direttamente nella directory contenente la versione
iniziale dei propri sorgenti nel qual caso il comando richiede come ulteriore
argomento la directory o la URL con la quale indicare il \textit{repository}
da usare. Alternativamente si può passare come primo argomento il pathname
della directory da importare, seguito dall'indicazione della URL del
\textit{repository}.

Si tenga presente che l'operazione di importazione inserisce sul
\textit{repository} il contenuto completo della directory indicata, compresi
eventuali file nascosti e sotto-directory. È anche possibile eseguire
l'importazione di più directory da inserire in diverse sezioni del
\textit{repository}, ma un tal caso ciascuna importazione sarà vista con una
diversa \textit{release}. Ad ogni operazione di modifica del
\textit{repository} viene infatti assegnato un numero progressivo che consente
di identificarne la storia delle modifiche e riportarsi ad un dato punto della
stessa in ogni momento successivo.\footnote{a differenza di CVS Subversion non
  assegna un numero di versione progressivo distinto ad ogni file, ma un
  numero di \textit{release} progressivo ad ogni cambiamento globale del
  repository, pertanto non esiste il concetto di versione di un singolo file,
  quanto di stato di tutto il \textit{repository} ad un dato momento, è
  comunque possibile richiedere in maniera indipendente la versione di ogni
  singolo file a qualunque \textit{release} si desideri.}

Una volta eseguita l'importazione di una versione iniziale è d'uopo cancellare
la directory originale e ripartire dal progetto appena creato. L'operazione di
recuperare ex-novo di tutti i file che fanno parte di un progetto, chiamata
usualmente \textit{checkout}, viene eseguita con il
comando:\footnote{alternativamente si può usare l'abbreviazione \texttt{svn
    co}.}
\begin{verbatim}
svn checkout URL [/pathname]
\end{verbatim}
che creerà nella directory corrente una directory corrispondente al nome
specificato in coda alla URL passata come argomento, scaricando l'ultima
versione dei file archiviati sul repository; alternativamente si può
specificare come ulteriore argomento la directory su cui scaricare i file.

Sia in caso di \texttt{import} che di \texttt{checkout} è sempre possibile
operare su una qualunque sotto cartella contenuta all'interno di un
\textit{repository}, ignorando totalmente quello che sta al di sopra, basterà
indicare in sede di importazione o di estrazione iniziale un pathname o una
URL che identifichi quella parte del progetto.

Se quando si effettua lo scaricamento non si vuole usare la versione più
aggiornata, ma una versione precedente si può usare l'opzione \texttt{-r}
seguita da un numero che scaricherà esattamente quella \textit{release},
alternativamente al posto del numero si può indicare una data, e verrà presa
la \textit{release} più prossima a quella data. 

A differenza di CVS Subversion non supporta l'uso di \textsl{etichette}
associate ad una certa versione del proprio progetto, per questo è invalso
l'uso di strutturare il repository secondo lo schema illustrato inizialmente;
è infatti molto semplice (e non comporta nessun tipo di aggravio) creare delle
copie complete di una qualunque parte del \textit{repository} su un'altra
parte dello stesso, per cui se si è eseguito lo sviluppo sulla cartella
\texttt{trunk} sarà possibile creare banalmente una versione con etichetta
\texttt{label} (o quel che si preferisce) semplicemente con una copia eseguita
con:
\begin{verbatim}
svn cp trunk tags/label
\end{verbatim}

Il risultato di questo comando è la creazione della nuova cartella
\texttt{label} sotto \texttt{tags}, che sarà assolutamente identica, nel
contenuto (e nella sua storia) a quanto presente in \texttt{trunk} al momento
dell'esecuzione del comando. In questo modo, una volta salvate le
modifiche,\footnote{la copia viene eseguita localmente verrà creata anche sul
  \textit{repository} solo dopo un \textit{commit}.} si potrà ottenere la
versione \texttt{label} del proprio progetto semplicemente eseguendo un
\textit{checkout} di \texttt{tags/label} in un'altra
directory.\footnote{ovviamente una volta presa la suddetta versione si deve
  aver cura di non eseguire nessuna modifica a partire dalla stessa, per
  questo se si deve modificare una versione etichettata si usa
  \texttt{branches}.}

Una volta creata la propria copia locale dei programmi, è possibile lavorare
su di essi ponendosi nella relativa directory, e apportare tutte le modifiche
che si vogliono ai file ivi presenti; due comandi permettono inoltre di
schedulare la rimozione o l'aggiunta di file al
\textit{repository}:\footnote{a differenza di CVS si possono aggiungere e
  rimuovere, ed anche spostare con \texttt{svn mv}, sia file che directory.}
\begin{verbatim}
svn add file1.c
svn remove file2.c
\end{verbatim}
ma niente viene modificato sul \textit{repository} fintanto che non viene
eseguito il cosiddetto \textit{commit} delle modifiche, vale a dire fintanto
che non viene dato il comando:\footnote{in genere anche questo viene
  abbreviato, con \texttt{svn ci}.}
\begin{verbatim}
svn commit [file]
\end{verbatim}
ed è possibile eseguire il \textit{commit} delle modifiche per un singolo
file, indicandolo come ulteriore argomento, mentre se non si indica nulla
verranno inviate tutte le modifiche presenti.

Si tenga presente però che il \textit{commit} non verrà eseguito se nel
frattempo i file del \textit{repository} sono stati modificati; in questo caso
\texttt{svn} rileverà la presenza di differenze fra la propria
\textit{release} e quella del \textit{repository} e chiederà che si effettui
preventivamente un aggiornamento. Questa è una delle operazioni di base di
Subversion, che in genere si compie tutte le volte che si inizia a lavorare,
il comando che la esegue è:
\begin{verbatim}
svn update
\end{verbatim}

Questo comando opera a partire dalla directory in cui viene eseguito e ne
aggiorna il contenuto (compreso quello di eventuali sotto-directory) alla
versione presente, scaricando le ultime versioni dei file esistenti o nuovi
file o directory aggiunte, cancellando eventuali file e directory rimossi dal
\textit{repository}.  Esso inoltre esso cerca, in caso di presenza di
modifiche eseguite in maniera indipendente sulla propria copia locale, di
eseguire un \textsl{raccordo} (il cosiddetto \textit{merging}) delle stesse
con quelle presenti sulla versione del \textit{repository}.

Fintanto che sono state modificate parti indipendenti di un file di testo in
genere il processo di \textit{merging} ha successo e le modifiche vengono
incorporate automaticamente in conseguenza dell'aggiornamento, ma quando le
modifiche attengono alla stessa parte di un file nel ci si troverà di fronte
ad un conflitto ed a quel punto sarà richiesto al ``committente'' di
intervenire manualmente sui file per i quali sono stati rilevati i conflitti
per risolverli.

Per aiutare il committente nel suo compito quando l'operazione di
aggiornamento fallisce nel raccordo delle modifiche lascia sezioni di codice
in conflitto opportunamente marcate e separate fra loro come nell'esempio
seguente:
\begin{verbatim}
<<<<<<< .mine
        $(CC) $(CFLAGS) -o pamacq pamacq.c -lm
=======
        $(CC) $(CFLAGS) -o pamacq pamacq.c

>>>>>>> r.122
\end{verbatim}

In questo caso si c'è stata una modifica sul file (mostrata nella parte
superiore) incompatibile con quella fatta nel \textit{repository} (mostrata
nella parte inferiore). Prima di eseguire un \textit{commit} occorrerà
pertanto integrare le modifiche e salvare nuovamente il file rimuovendo i
marcatori, inoltre prima che il \textit{commit} ritorni possibile si dovrà
esplicitare la risoluzione del conflitto con il comando:
\begin{verbatim}
svn resolved file
\end{verbatim}

\begin{table}[htb]
  \centering
  \footnotesize
  \begin{tabular}[c]{|c|l|}
    \hline
    \textbf{Flag} & \textbf{Significato} \\
    \hline
    \hline
    \texttt{?}& File sconosciuto.\\
    \texttt{M}& File modificato localmente.\\
    \texttt{A}& File aggiunto.\\
    \texttt{C}& File con conflitto.\\
%    \const{G}& File con modifiche integrate.\\
%    \const{}& .\\
    \hline
  \end{tabular}
  \caption{Caratteri associati ai vari stati dei file.}
  \label{tab:svn_status}
\end{table}


Infine per capire la situazione della propria copia locale si può utilizzare
il comando \texttt{svn status} che confronta i file presenti nella directory
locale rispetto alla ultima versione scaricata dal \textit{repository} e per
tutti quelli che non corrispondono stampa a schermo delle informazioni di
stato nella forma di un carattere seguito dal nome del file, secondo quanto
illustrato in tab.~\ref{tab:svn_status}.




% LocalWords:  make Makefile makefile shell FC FLAGS CFLAGS pathname CERN SCM
% LocalWords:  OBJ commondef readfile FFLAGS Cuncurrent Version System CVS home
% LocalWords:  adidsp geometry muonacq andaranno CVSROOT startup cvs ssh rsh to
% LocalWords:  project tag Root commit update merging Locally Modified Added co
% LocalWords:  Needs Patch Merge log rtag checkout module altrobersaglio Source
% LocalWords:   Control Subversion repository DAV svnserve URL svn subversion
% LocalWords:  client sez svnadmin Truelite Srl trunk tags branches tab add rm
% LocalWords:  remove copy cp move resolved label



%%% Local Variables: 
%%% mode: latex
%%% TeX-master: "gapil"
%%% End: 
